\subsection{Non-local form factors and anomaly}
  \label{subsec:nonlocal-anomaly}

  The finite part of the determinant generates non-local contributions.
  In covariant form, these can be expressed through curvature-dependent
  form factors of the type
  \begin{equation}
    R \,\mathcal{F}\!\left(\frac{\Box}{\mu^2}\right) R ,
  \end{equation}
  and similar structures, as developed by
  Barvinsky and Vilkovisky~\cite{BarvinskyVilkovisky1985,BarvinskyVilkovisky1990}.
  Their metric variation yields non-local tensors
  $T^{\mathrm{non\mbox{-}loc}}_{\mu\nu}$.

  In addition, renormalization induces a conformal trace anomaly.
  For a minimally coupled scalar field in four dimensions,
  the trace of the anomaly satisfies~\cite{BirrellDavies1982,Vassilevich2003}
  \begin{equation}
    g^{\mu\nu} A_{\mu\nu}
    =
    \frac{1}{(4\pi)^2}
    \frac{1}{360}
    \left(
      R_{\mu\nu\rho\sigma} R^{\mu\nu\rho\sigma}
      -
      R_{\mu\nu} R^{\mu\nu}
      +
      \frac{1}{3} \Box R
    \right).
  \end{equation}

  The local tensor structure of $A_{\mu\nu}$ depends on the
  renormalization scheme, but its trace is fixed by this universal expression.
  The anomaly contributes at four-derivative order and remains subleading
  in the weak-curvature infrared regime $R\,\ell_\chi^2 \ll 1$.
